% =====================================================================
%  符号说明(独立章节,编号随正文顺序)
%  由队长(Claude)统一维护,全队以此为准。
%  建模手写公式、代码手写代码、论文手排版都用同一套符号。
%  新符号出现时,更新此表并推送,大家在群里同步。
%
%  排版说明:符号较多,切分成两个表、分两页连续排放。
%  用 \clearpage 保证每表独占一页(浮动体队列清空,不产生空白页)。
% =====================================================================
\clearpage
\section{符号说明}

% ---------- 表一:传热与流动物理量 ----------
\begin{table}[h]
\centering
\caption{主要符号说明(一):传热与流动物理量}
\begin{tabular}{lll}
\toprule
符号 & 含义 & 单位 \\
\midrule
$\alpha$   & 针肋宽度比(针肋直径/通道宽度) & — \\
$\beta$    & 歧管深高比(歧管层/微通道层深高比) & — \\
$n$        & 单个歧管单元内沿流向的针肋排数 & — \\
$R_{th}$   & 总热阻 & K/W \\
$\Delta P$ & 压降 & Pa \\
$\Delta T$ & 温度非均匀性(最高-最低温度差) & K \\
$\sigma_T$ & 温度非均匀性(底面温度标准差,备选) & K \\
$T_{max}$  & 芯片表面最高温度 & K \\
$T_{min}$  & 芯片表面最低温度 & K \\
$T_{in}$   & 冷却液入口温度(293 K) & K \\
$T_f$      & 流体温度 & K \\
$T_s$      & 固体温度 & K \\
$P_{in}$   & 入口压力 & Pa \\
$P_{out}$  & 出口压力 & Pa \\
$Q$        & 芯片总发热功率 & W \\
$\vec{v}$  & 速度矢量 & m/s \\
$h$        & 对流换热系数 & W/(m$^2$$\cdot$K) \\
$A$        & 换热面积 & m$^2$ \\
$D_h$      & 通道水力直径 & m \\
$u$        & 流速 & m/s \\
$\rho$     & 冷却液密度 & kg/m$^3$ \\
$C_p$      & 冷却液比热容 & J/(kg$\cdot$K) \\
$\mu$      & 冷却液动力黏度 & Pa$\cdot$s \\
$k_f$      & 流体导热系数 & W/(m$\cdot$K) \\
$k_s$      & 固体(衬底)导热系数 & W/(m$\cdot$K) \\
$f$        & 沿程阻力系数 & — \\
$Re$       & 雷诺数 & — \\
\bottomrule
\end{tabular}
\end{table}

\clearpage

% ---------- 表二:代理模型量 ----------
\begin{table}[h]
\centering
\caption{主要符号说明(二):代理模型量}
\begin{tabular}{lll}
\toprule
符号 & 含义 & 单位 \\
\midrule
$\mathbf{X}$    & 结构参数输入向量 $[x_1, x_2, x_3]^T$ & — \\
$\mathbf{Y}$    & 性能指标输出向量 $[y_1, y_2, y_3]^T$ & — \\
$\tilde{x}_i$   & 标准化后的第 $i$ 个输入特征 & — \\
$\mu_i$         & 第 $i$ 个特征在训练集上的均值 & — \\
$\sigma_i$      & 第 $i$ 个特征在训练集上的标准差 & — \\
$f_j(\cdot)$    & 第 $j$ 项性能指标的潜在函数 & — \\
$\mathcal{GP}$  & 高斯过程 & — \\
$k(\cdot, \cdot)$ & 核函数 & — \\
$\sigma_f^2$    & 信号方差(核函数幅值) & — \\
$l_i$           & 第 $i$ 个结构参数的各向异性长度尺度 & — \\
$\sigma_n^2$    & 观测噪声方差 & — \\
$\delta(\cdot, \cdot)$ & Kronecker delta 函数 & — \\
$\boldsymbol{\theta}$ & 超参数集合 $[\sigma_f, l_1, l_2, l_3, \sigma_n]$ & — \\
$\mathbf{K}$    & 训练集核矩阵 & — \\
$\mathbf{k}$    & 训练集与新样本间的协方差向量 & — \\
$\mathbf{I}$    & 单位矩阵 & — \\
$\mathbf{y}_{train}$ & 训练集输出向量 & — \\
$R^2$           & 决定系数 & — \\
$RMSE$          & 均方根误差 & — \\
$MAPE$          & 平均绝对百分比误差 & — \\
$\hat{y}_k$     & 第 $k$ 个样本的模型预测值 & — \\
$\bar{y}$       & 观测均值 & — \\
$N$             & 样本数 & — \\
\bottomrule
\end{tabular}
\end{table}

% 说明:以上为预填常用符号,比赛过程中随建模进度增删。
% 新增符号时:告知队长 → 队长更新本表 → 推送,全队同步。
